\section{Experimental Setup}

Our experiments validate a core hypothesis: zero-shot evaluation reveals architectural compatibility before parameter-efficient fine-tuning investment. We design experiments to answer specific research questions connecting to our central claims.

\subsection{Research Questions}

\textbf{RQ1: Does Mamba-130M demonstrate task-dependent zero-shot compatibility?}
We hypothesize that causal state-space models succeed on generation-aligned classification tasks (e.g., sentiment analysis) but fail on tasks requiring bidirectional reasoning (e.g., paraphrase detection). Success is defined as above-random baseline accuracy; failure is below-random accuracy indicating systematic architectural bias.

\textbf{RQ2: Is below-random performance statistically significant or sampling noise?}
We test whether observed failures reflect genuine architectural constraints or insufficient sample sizes. Binomial significance tests distinguish systematic failure patterns from random variation.

\textbf{RQ3: Does zero-shot performance correlate with task structure?}
We predict that task requirements (symmetric comparison vs. unidirectional encoding) determine success/failure patterns independent of task difficulty or domain.

\subsection{Tasks and Datasets}

We evaluate on three GLUE benchmark tasks \cite{wang2018glue} representing diverse reasoning requirements:

\textbf{QQP (Quora Question Pairs):} Binary paraphrase detection over question pairs. Requires symmetric comparison---both questions must inform similarity judgment equally. Random baseline: 50\%. We sample 100 validation examples for rapid evaluation.

\textbf{MNLI (Multi-Genre Natural Language Inference):} Three-way classification (entailment, neutral, contradiction) over premise-hypothesis pairs. Like QQP, requires bidirectional reasoning to compare premise and hypothesis symmetrically. Random baseline: 33.3\%. We evaluate on 100 matched validation examples.

\textbf{SST-2 (Stanford Sentiment Treebank):} Binary sentiment classification over single sentences. Aligns with causal language modeling---left-to-right encoding aggregates context sufficient for classification. Random baseline: 50\%. We evaluate on 100 validation examples.

These tasks form a minimal test suite isolating architectural effects: QQP and MNLI probe bidirectional reasoning (predicted failure for causal SSMs), while SST-2 tests generation-aligned classification (predicted success).

\begin{table}[h]
\centering
\small
\begin{tabular}{lcccc}
\toprule
Task & Type & \# Classes & Baseline & Requirement \\
\midrule
QQP & Paraphrase & 2 & 50\% & Symmetric comparison \\
MNLI & NLI & 3 & 33.3\% & Bidirectional reasoning \\
SST-2 & Sentiment & 2 & 50\% & Unidirectional encoding \\
\bottomrule
\end{tabular}
\caption{Task overview showing architectural requirements.}
\label{tab:tasks}
\end{table}

\subsection{Model Configuration}

\textbf{Mamba-130M:} We use the pretrained checkpoint \texttt{state-spaces/mamba-130m-hf} from HuggingFace. The model employs selective state-space mechanisms with 24 layers processing 130M total parameters. Pretrained on standard language modeling corpora using causal autoregressive objectives.

\textbf{Classification heads:} For each task, we attach a randomly initialized linear layer mapping final hidden states to class logits. No head pretraining or fine-tuning is performed---true zero-shot evaluation.

\textbf{Inference protocol:} Deterministic evaluation with seed 42, greedy decoding, no prompt engineering or few-shot examples. We measure raw zero-shot capability without task-specific optimization.

\subsection{Evaluation Metrics}

\textbf{Primary metric:} Accuracy (percentage of correct predictions). Compared against random baseline to determine architectural compatibility.

\textbf{Statistical significance:} Binomial tests assess whether observed accuracy deviates significantly from random performance. For binary tasks (QQP, SST-2), random baseline is 50\%; for MNLI, 33.3\%. We use $\alpha = 0.05$ significance threshold.

\textbf{Interpretation thresholds:}
\begin{itemize}
\item Accuracy $\geq$ random baseline + 5pp: Model demonstrates task capability (architecture compatible)
\item Accuracy $<$ random baseline: Systematic failure indicating architectural incompatibility
\item Accuracy within $\pm$5pp of baseline: Inconclusive (insufficient signal)
\end{itemize}

\subsection{Implementation Details}

All experiments run on a single NVIDIA A100 GPU using HuggingFace Transformers library (version 4.30.0). For each task:

\begin{enumerate}
\item Load pretrained Mamba-130M checkpoint
\item Initialize task-specific classification head
\item Perform forward pass on 100 validation examples (no gradient updates)
\item Compute accuracy and binomial significance
\end{enumerate}

Total evaluation time: $<$5 minutes per task. This efficiency demonstrates the practicality of zero-shot architectural compatibility testing as a PEFT prerequisite.

\subsection{Experimental Validity}

\textbf{Sample size justification:} 100 examples per task provides sufficient statistical power to detect below-random performance. For QQP, observing 38/100 correct (vs. expected 50/100 under null hypothesis) yields $p = 0.003$, confirming robust significance.

\textbf{Checkpoint selection:} We use the standard publicly available Mamba checkpoint rather than selecting from multiple variants, preventing cherry-picking bias. Results reflect base architectural capabilities, not checkpoint-specific quirks.

\textbf{No hyperparameter tuning:} Zero-shot evaluation removes confounds from hyperparameter optimization. Observed failures cannot be attributed to suboptimal learning rates or batch sizes---the architecture either supports the task or it doesn't.

\textbf{Reproducibility:} All experiments use deterministic seeds and publicly available datasets/checkpoints. Full implementation code released for verification.

\subsection{Baseline Comparisons}

We establish architectural compatibility through comparison to random baselines rather than other models. This design choice isolates architectural capability: below-random accuracy definitively indicates incompatibility independent of model size, pretraining data, or training duration. Comparisons to GPT-2 or other baselines would confound architectural effects with checkpoint quality differences.

Future work could extend this protocol to comprehensive baseline comparisons (GPT-2 LoRA, RoBERTa, etc.), but our focus is on validating the zero-shot gate principle, not benchmarking competitive performance.
